\documentclass[a4paper,12pt]{article}

% --- preprost slovenski LaTeX stil ---
\usepackage[utf8]{inputenc}
\usepackage[T1]{fontenc}
\usepackage[slovene]{babel}
\usepackage{lmodern}
\usepackage[a4paper,margin=2.5cm]{geometry}
\usepackage{amsmath}
\usepackage{amssymb}
\usepackage{enumitem}
\usepackage{amsthm}

\newtheoremstyle{naloga}
  {12pt}{6pt}{}{}{\bfseries}{.}{.5em}{}
\theoremstyle{naloga}
\newtheorem{naloga}{Naloga}[section]

\setlist[enumerate,1]{label=\alph*), leftmargin=1.5em}

\title{Geometrijska telesa preverjanje}
\author{}
\date{}

\begin{document}
\maketitle

\section{Test 2.0}

\begin{naloga}[4+5 t.]
Pravilni šesterokotni prizmi je osnovni rob $a=6$ cm, višina prizme pa je enaka premeru vrisanega kroga osnovne ploskve.
\begin{enumerate}
\item Izračunaj polmer vrisanega kroga in višino prizme.
\item Izračunaj prostornino in površino prizme.
\end{enumerate}
\end{naloga}

\begin{naloga}[3+4+4 t.]
Pravilni štiristrani piramidi meri diagonala osnovne ploskve $d=12$ cm, stranski rob pa $b=10$ cm.
\begin{enumerate}
\item Izračunaj rob osnovne ploskve in višino piramide.
\item Izračunaj prostornino piramide.
\item Izračunaj kot med stranskim robom in osnovno ploskvijo (do kotne minute natančno).
\end{enumerate}
\end{naloga}

\begin{naloga}[4+4 t.]
V stožec s polmerom osnovne ploskve $r=6$ cm in višino $v=8$ cm vrišemo kroglo (dotika se plašča in osnovne ploskve).
\begin{enumerate}
\item Izračunaj polmer vrisane krogle.
\item V kakšnem razmerju sta prostornini krogle in stožca?
\end{enumerate}
\end{naloga}

\begin{naloga}[4+4 t.]
Pravokotni trapez s pravokotnim krakom $12$ cm ter osnovnicama $9$ cm in $4$ cm zavrtimo okoli daljšega (pravokotnega) kraka za polni kot.
\begin{enumerate}
\item Izračunaj prostornino nastale vrtenine.
\item Izračunaj površino nastale vrtenine.
\end{enumerate}
\end{naloga}

\section{Test 2.1}

\begin{naloga}[4+5 t.]
Pravilni tristrani prizmi je rob osnovne ploskve $a=8$ cm. Diagonala stranske ploskve oklepa z osnovno ploskvijo kot $60^{\circ}$.
\begin{enumerate}
\item Izračunaj višino prizme.
\item Izračunaj prostornino in površino prizme.
\end{enumerate}
\end{naloga}

\begin{naloga}[3+5 t.]
Pravilno štiristrano piramido z osnovnim robom $a=12$ cm in višino $v=9$ cm prerežemo z ravnino, vzporedno z osnovno ploskvijo, na polovici višine.
\begin{enumerate}
\item Izračunaj prostornino cele piramide.
\item V kakšnem razmerju sta prostornini zgornje (male) piramide in prisekane piramide (spodnjega dela)?
\end{enumerate}
\end{naloga}

\begin{naloga}[4+4 t.]
Plašč stožca razgrnemo v krožni izsek s polmerom (stranico stožca) $s=15$ cm in središčnim kotom $216^{\circ}$.
\begin{enumerate}
\item Izračunaj polmer osnovne ploskve stožca in višino.
\item Izračunaj prostornino in površino stožca.
\end{enumerate}
\end{naloga}

\begin{naloga}[4+3 t.]
Krogli s polmerom $R=6$ cm očrtamo pokončni valj, tako da se krogla dotika obeh osnovnih ploskev in plašča valja.
\begin{enumerate}
\item Izračunaj prostornino in površino valja.
\item Kolikšen del prostornine valja zavzema krogla?
\end{enumerate}
\end{naloga}

\section{Test 2.2}

\begin{naloga}[3+4+4 t.]
Kvader ima robova osnovne ploskve $9$ cm in $12$ cm, prostorninska diagonala pa meri $17$ cm.
\begin{enumerate}
\item Izračunaj višino kvadra.
\item Izračunaj prostornino in površino kvadra.
\item Izračunaj kot med prostorninsko diagonalo in osnovno ploskvijo (do kotne minute natančno).
\end{enumerate}
\end{naloga}

\begin{naloga}[4+4 t.]
Pravilna štiristrana piramida ima vse robove enake dolžine $a=6$ cm (enakorobna piramida).
\begin{enumerate}
\item Izračunaj višino in prostornino piramide.
\item Izračunaj kot med stransko ploskvijo in osnovno ploskvijo (do kotne minute natančno).
\end{enumerate}
\end{naloga}

\begin{naloga}[4+4 t.]
Stožcu s polmerom osnovne ploskve $r=5$ cm in višino $v=12$ cm na osnovno ploskev prilepimo polkroglo z enakim polmerom.
\begin{enumerate}
\item Izračunaj prostornino sestavljenega telesa.
\item Izračunaj površino sestavljenega telesa (brez skupne ploskve).
\end{enumerate}
\end{naloga}

\begin{naloga}[4+4 t.]
Enakokraki trikotnik z osnovnico $16$ cm in krakoma $17$ cm zavrtimo okoli osnovnice za polni kot.
\begin{enumerate}
\item Izračunaj prostornino nastalega telesa.
\item Izračunaj površino nastalega telesa.
\end{enumerate}
\end{naloga}

\section{Test 2.3}

\begin{naloga}[4+3 t.]
Pravilni šesterokotni prizmi z osnovnim robom $a=4$ cm in višino $v=10$ cm je vrisan valj (dotika se vseh šestih stranskih ploskev).
\begin{enumerate}
\item Izračunaj polmer vrisanega valja in njegovo prostornino.
\item Koliko odstotkov prostornine prizme zavzema valj?
\end{enumerate}
\end{naloga}

\begin{naloga}[3+3+4 t.]
Pravilni tristrani piramidi je rob osnovne ploskve $a=6$ cm, stranski rob pa $b=4$ cm.
\begin{enumerate}
\item Izračunaj višino piramide.
\item Izračunaj prostornino piramide.
\item Izračunaj kot med stranskim robom in osnovno ploskvijo.
\end{enumerate}
\end{naloga}

\begin{naloga}[3+4+3 t.]
Osni presek stožca je enakostranični trikotnik s stranico $10$ cm.
\begin{enumerate}
\item Izračunaj polmer in višino stožca.
\item Izračunaj prostornino in površino stožca.
\item Izračunaj kot med stranico stožca in osnovno ploskvijo.
\end{enumerate}
\end{naloga}

\begin{naloga}[4+4 t.]
Lik, sestavljen iz pravokotnika s stranicama $8$ cm in $4$ cm ter polkroga s polmerom $4$ cm, ki je pripet na krajšo stranico pravokotnika, zavrtimo okoli daljše stranice (osi simetrije) za polni kot.
\begin{enumerate}
\item Izračunaj prostornino nastalega telesa.
\item Izračunaj površino nastalega telesa.
\end{enumerate}
\end{naloga}

\section{Test 2.4}

\begin{naloga}[3+4 t.]
Štiristrani pokončni prizmi je za osnovno ploskev pravokotni trapez z vzporednima stranicama $9$ cm in $6$ cm, pravokotnim krakom (višino trapeza) $4$ cm in poševnim krakom $5$ cm. Rob (višina) prizme meri $9$ cm.
\begin{enumerate}
\item Izračunaj ploščino osnovne ploskve.
\item Izračunaj prostornino in površino prizme.
\end{enumerate}
\end{naloga}

\begin{naloga}[4+4 t.]
Pravilna šesterokotna piramida ima osnovni rob $a=6$ cm in višino $v=4\sqrt{3}$ cm.
\begin{enumerate}
\item Izračunaj prostornino piramide.
\item Izračunaj površino piramide.
\end{enumerate}
\end{naloga}

\begin{naloga}[3+3+4 t.]
Stožec ima polmer osnovne ploskve $r=12$ cm in višino $v=16$ cm.
\begin{enumerate}
\item Izračunaj stranico stožca in njegovo površino.
\item Izračunaj prostornino stožca.
\item Izračunaj kot med stranico in višino (osjo) stožca, do kotne minute natančno.
\end{enumerate}
\end{naloga}

\begin{naloga}[3+3+4 t.]
Krogla in valj imata enako prostornino. Polmer krogle je $R=6$ cm, polmer valja pa je enak polmeru krogle.
\begin{enumerate}
\item Izračunaj prostornino krogle.
\item Izračunaj višino valja.
\item V kakšnem razmerju sta površini krogle in valja (plašč in obe osnovni ploskvi)?
\end{enumerate}
\end{naloga}

\clearpage
\section*{Rešitve}
\addcontentsline{toc}{section}{Rešitve}

\subsection*{Test 2.0}
\textbf{Naloga 1.1.} a) $r=3\sqrt3\ \text{cm}\approx5{,}20$ cm, $v=6\sqrt3\ \text{cm}\approx10{,}39$ cm.\quad
b) $V=972\ \text{cm}^3$, $S=324\sqrt3\ \text{cm}^2\approx561{,}18\ \text{cm}^2$.

\textbf{Naloga 1.2.} a) $a=6\sqrt2\ \text{cm}\approx8{,}49$ cm, $v=8$ cm.\quad
b) $V=192\ \text{cm}^3$.\quad
c) $\varphi\approx53^{\circ}8'$.

\textbf{Naloga 1.3.} a) $r_{\text{kr}}=3$ cm.\quad
b) $V_{\text{krogla}}:V_{\text{stožec}}=3:8$.

\textbf{Naloga 1.4.} a) $V=532\pi\ \text{cm}^3\approx1671{,}33\ \text{cm}^3$.\quad
b) $S=266\pi\ \text{cm}^2\approx835{,}66\ \text{cm}^2$.

\subsection*{Test 2.1}
\textbf{Naloga 2.1.} a) $v=8\sqrt3\ \text{cm}\approx13{,}86$ cm.\quad
b) $V=384\ \text{cm}^3$, $S=224\sqrt3\ \text{cm}^2\approx387{,}98\ \text{cm}^2$.

\textbf{Naloga 2.2.} a) $V_{\text{cela}}=432\ \text{cm}^3$ ($V_{\text{mala}}=54\ \text{cm}^3$, $V_{\text{prisekana}}=378\ \text{cm}^3$).\quad
b) razmerje $=1:7$.

\textbf{Naloga 2.3.} a) $r=9$ cm, $v=12$ cm.\quad
b) $V=324\pi\ \text{cm}^3\approx1017{,}88\ \text{cm}^3$, $S=216\pi\ \text{cm}^2\approx678{,}58\ \text{cm}^2$.

\textbf{Naloga 2.4.} a) $V_{\text{valj}}=432\pi\ \text{cm}^3\approx1357{,}17\ \text{cm}^3$, $S_{\text{valj}}=216\pi\ \text{cm}^2\approx678{,}58\ \text{cm}^2$.\quad
b) krogla zavzema $\dfrac{2}{3}$ prostornine valja ($\approx66{,}7\,\%$).

\subsection*{Test 2.2}
\textbf{Naloga 3.1.} a) $v=8$ cm (diagonala osnovne ploskve $=15$ cm).\quad
b) $V=864\ \text{cm}^3$, $S=552\ \text{cm}^2$.\quad
c) $\varphi\approx28^{\circ}4'$.

\textbf{Naloga 3.2.} a) $v=3\sqrt2\ \text{cm}\approx4{,}24$ cm, $V=9\sqrt6\ \text{cm}^3\approx22{,}05\ \text{cm}^3$.\quad
b) $\varphi\approx54^{\circ}44'$.

\textbf{Naloga 3.3.} a) $V=\dfrac{550\pi}{3}\ \text{cm}^3\approx575{,}96\ \text{cm}^3$ ($V_{\text{stožec}}=100\pi$, $V_{\text{polkrogla}}=\tfrac{250\pi}{3}$).\quad
b) $S=115\pi\ \text{cm}^2\approx361{,}28\ \text{cm}^2$ ($S_{\text{plašč}}=65\pi$, $S_{\text{polkrogla}}=50\pi$).

\textbf{Naloga 3.4.} a) $V=1200\pi\ \text{cm}^3\approx3769{,}91\ \text{cm}^3$.\quad
b) $S=510\pi\ \text{cm}^2\approx1602{,}21\ \text{cm}^2$.

\subsection*{Test 2.3}
\textbf{Naloga 4.1.} a) $r=2\sqrt3\ \text{cm}\approx3{,}46$ cm, $V_{\text{valj}}=120\pi\ \text{cm}^3\approx376{,}99\ \text{cm}^3$.\quad
b) $V_{\text{prizma}}=240\sqrt3\ \text{cm}^3\approx415{,}69\ \text{cm}^3$; valj zavzema $\approx90{,}69\,\%$ prostornine prizme.

\textbf{Naloga 4.2.} a) $v=2$ cm.\quad
b) $V=6\sqrt3\ \text{cm}^3\approx10{,}39\ \text{cm}^3$.\quad
c) $\varphi=30^{\circ}$.

\textbf{Naloga 4.3.} a) $r=5$ cm, $v=5\sqrt3\ \text{cm}\approx8{,}66$ cm.\quad
b) $V=\dfrac{125\sqrt3}{3}\pi\ \text{cm}^3\approx226{,}72\ \text{cm}^3$, $S=75\pi\ \text{cm}^2\approx235{,}62\ \text{cm}^2$.\quad
c) $\varphi=60^{\circ}$.

\textbf{Naloga 4.4.} a) $V=\dfrac{512\pi}{3}\ \text{cm}^3\approx536{,}17\ \text{cm}^3$ ($V_{\text{valj}}=128\pi$, $V_{\text{polkrogla}}=\tfrac{128\pi}{3}$).\quad
b) $S=112\pi\ \text{cm}^2\approx351{,}86\ \text{cm}^2$.

\subsection*{Test 2.4}
\textbf{Naloga 5.1.} a) ploščina osnovne ploskve $=30\ \text{cm}^2$.\quad
b) $V=270\ \text{cm}^3$, $S=276\ \text{cm}^2$.

\textbf{Naloga 5.2.} a) $V=216\ \text{cm}^3$.\quad
b) $S=144\sqrt3\ \text{cm}^2\approx249{,}42\ \text{cm}^2$.

\textbf{Naloga 5.3.} a) $s=20$ cm, $S=384\pi\ \text{cm}^2\approx1206{,}37\ \text{cm}^2$.\quad
b) $V=768\pi\ \text{cm}^3\approx2412{,}74\ \text{cm}^3$.\quad
c) $\varphi\approx36^{\circ}52'$.

\textbf{Naloga 5.4.} a) $V_{\text{krogla}}=288\pi\ \text{cm}^3\approx904{,}78\ \text{cm}^3$.\quad
b) $v_{\text{valj}}=8$ cm.\quad
c) $S_{\text{krogla}}:S_{\text{valj}}=6:7$ ($S_{\text{krogla}}=144\pi$, $S_{\text{valj}}=168\pi$).

\end{document}
